\documentclass[12pt]{article}
\usepackage[utf8]{inputenc}
\usepackage[T1]{fontenc}
\usepackage[margin=1in]{geometry}
\usepackage{setspace}
\usepackage{xcolor}
\usepackage[colorlinks=true, urlcolor=blue!60!black, linkcolor=blue!60!black]{hyperref}
\setlength{\parindent}{2em}
\setlength{\parskip}{1.25em}
\renewcommand{\baselinestretch}{1.0}
\begin{document}
\begin{center}
{\LARGE \textit{Seminole Tribe of Florida v. Florida}, 517 U.S. 44 (1996)\par}
\vspace{0.8em}
{\large Case Brief\par}
\end{center}
\vspace{1em}

\section{Memory Jogger}

Overruling \textit{Pennsylvania v. Union Gas Co.}, the Court held that Congress cannot use its Article I powers to abrogate state sovereign immunity under the Eleventh Amendment---leaving Section 5 of the Fourteenth Amendment as the primary constitutional vehicle for abrogation---and that \textit{Ex parte Young} cannot be used to bypass an intricate statutory remedial scheme (\href{https://www.courtlistener.com/opinion/118011/seminole-tribe-of-florida-v-florida/}{CourtListener}; \href{https://www.law.cornell.edu/supct/html/94-12.ZO.html}{Cornell LII}).

\section{Detailed Case Facts}

In 1988, acting pursuant to the Indian Commerce Clause (U.S. Const., Art. I, \S 8, cl. 3), Congress enacted the Indian Gaming Regulatory Act (IGRA), 102 Stat. 2475, 25 U.S.C. \S 2701 \textit{et seq.}, to establish a federal statutory framework for the operation and regulation of gaming on Indian lands (\href{https://www.law.cornell.edu/supct/html/94-12.ZO.html}{Cornell LII}). IGRA divided tribal gaming into three classes: Class I (social and traditional ceremonial games under exclusive tribal jurisdiction), Class II (bingo and certain nonbanking card games), and Class III (the most heavily regulated category, encompassing slot machines, casino games, banking card games, dog racing, and lotteries) (25 U.S.C. \S 2703(8)) (\href{https://www.law.cornell.edu/supct/html/94-12.ZO.html}{Cornell LII}).

Under 25 U.S.C. \S 2710(d)(1), Class III gaming is lawful only if it is (1) authorized by an approved tribal ordinance, (2) located in a State that permits such gaming for any purpose by any person or entity, and (3) conducted in conformance with a valid Tribal-State compact entered into under \S 2710(d)(3) and approved by the Secretary of the Interior (\href{https://www.law.cornell.edu/supct/html/94-12.ZO.html}{Cornell LII}). To prevent States from vetoing tribal gaming by refusing to negotiate, \S 2710(d)(3)(A) imposed an affirmative duty on a State, upon a tribe's request, to ``negotiate with the Indian tribe in good faith to enter into such a compact.''

Congress made that good-faith duty judicially enforceable through an intricate remedial scheme in 25 U.S.C. \S 2710(d)(7). Specifically, \S 2710(d)(7)(A)(i) vested federal district courts with jurisdiction over ``any cause of action initiated by an Indian tribe arising from the failure of a State to enter into negotiations . . . or to conduct such negotiations in good faith,'' exercisable 180 days after the tribe's request (\href{https://www.law.cornell.edu/supct/html/94-12.ZO.html}{Cornell LII}). If a tribe showed that no compact existed and the State failed to respond in good faith, the burden shifted to the State to prove good faith (\S 2710(d)(7)(B)(ii)). If the district court found a lack of good faith, it could order the State and tribe to conclude a compact within 60 days (\S 2710(d)(7)(B)(iii)). If they failed, each side would submit a ``last best offer'' to a court-appointed mediator, who would pick the offer best comporting with IGRA (\S 2710(d)(7)(B)(iv)--(v)). If the State still did not consent within 60 days, the mediator would notify the Secretary of the Interior, who would prescribe Class III gaming procedures for the tribal lands (\S 2710(d)(7)(B)(vii)) (\href{https://www.law.cornell.edu/supct/html/94-12.ZO.html}{Cornell LII}).

In September 1991, the Seminole Tribe of Florida sued the State of Florida and its Governor, Lawton Chiles, in federal district court under 25 U.S.C. \S 2710(d)(7)(A) and 28 U.S.C. \S\S 1331 and 1362, alleging that Florida and Governor Chiles had refused to enter into any negotiations to include certain Class III gaming activities in a Tribal-State compact, thereby violating \S 2710(d)(3)'s good-faith negotiation requirement (\href{https://www.law.cornell.edu/supct/html/94-12.ZO.html}{Cornell LII}).

\section{Procedural History}

Florida and Governor Chiles moved to dismiss the complaint on the ground that the suit violated the State's sovereign immunity under the Eleventh Amendment. The U.S. District Court for the Southern District of Florida denied the motion to dismiss (\textit{Seminole Tribe of Florida v. Florida}, 801 F. Supp. 655 (S.D. Fla. 1992)), holding that Congress had validly abrogated state sovereign immunity under the Indian Commerce Clause (\href{https://www.law.cornell.edu/supct/html/94-12.ZO.html}{Cornell LII}; \href{https://supreme.justia.com/cases/federal/us/517/44/}{Justia}). Florida and Governor Chiles took an immediate interlocutory appeal under the collateral order doctrine (\textit{Puerto Rico Aqueduct \& Sewer Auth. v. Metcalf \& Eddy, Inc.}, 506 U.S. 139 (1993)).

The U.S. Court of Appeals for the Eleventh Circuit consolidated the appeal with a parallel Alabama case (\textit{Poarch Band of Creek Indians v. Alabama}) and reversed (\textit{Seminole Tribe of Florida v. Florida}, 11 F.3d 1016 (11th Cir. 1994)) (\href{https://www.law.cornell.edu/supct/html/94-12.ZO.html}{Cornell LII}). Although the Eleventh Circuit agreed that Congress had unmistakably expressed its intent to abrogate sovereign immunity in \S 2710(d)(7), it held that the Indian Commerce Clause did not confer constitutional authority to abrogate Eleventh Amendment immunity and that \textit{Ex parte Young}, 209 U.S. 123 (1908), did not permit the Tribe to sue the Governor to compel good-faith negotiations. The Eleventh Circuit remanded with instructions to dismiss for lack of subject-matter jurisdiction (while severing the offending provisions to allow tribes direct recourse to the Secretary of the Interior, a separate point the Supreme Court did not review) (\href{https://www.law.cornell.edu/supct/html/94-12.ZO.html}{Cornell LII}).

The Supreme Court granted certiorari (513 U.S. 1125 (1995)) to decide (1) whether the Eleventh Amendment prevents Congress from authorizing tribal suits against unconsenting States under the Indian Commerce Clause, and (2) whether \textit{Ex parte Young} permits a suit against the Governor to enforce IGRA's good-faith bargaining duty. The case was argued on October 11, 1995, and decided on March 27, 1996 (\href{https://www.law.cornell.edu/supct/html/94-12.ZO.html}{Cornell LII}; \href{https://www.oyez.org/cases/1995/94-12}{Oyez}).

\section{Judicial Votes}

The Court decided the case by a 5--4 vote on all issues (\href{https://www.law.cornell.edu/supct/html/94-12.ZO.html}{Cornell LII}; \href{https://supreme.justia.com/cases/federal/us/517/44/}{Justia}):

\begin{itemize}
\item \textbf{Majority (5):} Chief Justice Rehnquist (author), joined by Justices O'Connor, Scalia, Kennedy, and Thomas.
\item \textbf{Dissents (4):}
\begin{itemize}
\item Justice Stevens filed a dissenting opinion (also endorsing Justice Souter's dissent).
\item Justice Souter filed a dissenting opinion, joined by Justices Ginsburg and Breyer.
\end{itemize}
\end{itemize}

\section{Holding}

\textbf{Judgment:} The judgment of the U.S. Court of Appeals for the Eleventh Circuit dismissing the Seminole Tribe's suit against both the State of Florida and Governor Lawton Chiles for lack of jurisdiction is affirmed (\href{https://www.law.cornell.edu/supct/html/94-12.ZO.html}{Cornell LII}).

\textbf{Rule Announced:} To abrogate state sovereign immunity, Congress must (1) unequivocally express its intent to do so in clear statutory text, and (2) act pursuant to a valid grant of constitutional authority. While IGRA satisfied the clear-statement requirement, the Court overruled \textit{Pennsylvania v. Union Gas Co.}, 491 U.S. 1 (1989), and held that the Eleventh Amendment restricts the federal judicial power under Article III such that \textit{Article I} powers (including the Indian Commerce Clause and Interstate Commerce Clause) cannot be used to abrogate the sovereign immunity of unconsenting States. By contrast, Congress retains authority to abrogate state sovereign immunity when validly exercising its enforcement power under \textit{Section 5 of the Fourteenth Amendment} (\textit{Fitzpatrick v. Bitzer}, 427 U.S. 445 (1976)), because the Fourteenth Amendment was ratified after the Eleventh Amendment and expressly altered the federal-state balance. Furthermore, where Congress has prescribed a detailed statutory remedial scheme for enforcing a federal statutory right against a State, federal courts may not allow plaintiffs to circumvent those limitations via a suit against a state officer under \textit{Ex parte Young} (\href{https://www.law.cornell.edu/supct/html/94-12.ZO.html}{Cornell LII}).

\textbf{Holding vs. Dicta:} The core \textit{holdings} are (1) Congress clearly intended to abrogate state immunity in 25 U.S.C. \S 2710(d)(7); (2) \textit{Pennsylvania v. Union Gas Co.} is overruled and neither the Indian Commerce Clause nor other Article I Commerce powers authorize abrogation of Eleventh Amendment immunity; and (3) \textit{Ex parte Young} is unavailable to enforce \S 2710(d)(3) against Florida's Governor in light of \S 2710(d)(7)'s intricate remedial scheme. Important \textit{dicta} and open questions include: (1) the Court's broad statement that \textit{no} Article I power can circumvent Article III's sovereign-immunity limits was later qualified when the Court carved out a narrow Bankruptcy Clause exception rooted in the ``plan of the Convention'' (\textit{Central Virginia Community College v. Katz}, 546 U.S. 356 (2006)) and an eminent-domain exception (\textit{PennEast Pipeline Co. v. New Jersey}, 594 U.S. 482 (2021)) (\href{https://texaslawreview.org/state-sovereign-immunity-after-the-revolution/}{Texas Law Review}); (2) the Court expressly declined to opine on the Eleventh Circuit's substitute administrative remedy allowing tribes to petition the Secretary of the Interior directly (footnote 18); and (3) the Court did not address Florida's Tenth Amendment anti-commandeering argument because it was not raised below (footnote 10) (\href{https://www.law.cornell.edu/supct/html/94-12.ZO.html}{Cornell LII}).

\section{Analysis of Opinions}

\subsection{Opinion of the Court --- Chief Justice Rehnquist}

Chief Justice Rehnquist began by emphasizing that although the Eleventh Amendment's literal text appears to restrict only diversity suits against States by citizens of other States or foreign subjects, the Court since \textit{Hans v. Louisiana}, 134 U.S. 1 (1890), has understood the Amendment to confirm a broader constitutional ``presupposition'': each State is a sovereign entity not amenable to private suit without its consent (\href{https://www.law.cornell.edu/supct/html/94-12.ZO.html}{Cornell LII}). Applying the two-part abrogation test from \textit{Green v. Mansour}, 474 U.S. 64 (1985), Rehnquist first held that IGRA's repeated references to ``the State'' as defendant in \S 2710(d)(7)(B) provided an ``unmistakably clear'' statement of congressional intent to abrogate (\href{https://www.law.cornell.edu/supct/html/94-12.ZO.html}{Cornell LII}).

Turning to constitutional power, Rehnquist rejected the Tribe's arguments that seeking prospective injunctive relief rather than money damages, or granting States new regulatory authority over Indian gaming, mattered to the Eleventh Amendment inquiry: sovereign immunity protects not only state treasuries, but also against ``the indignity of subjecting a State to the coercive process of judicial tribunals at the instance of private parties'' (\href{https://www.law.cornell.edu/supct/html/94-12.ZO.html}{Cornell LII}).

Rehnquist then contrasted the only two constitutional sources previously held to authorize abrogation: Section 5 of the Fourteenth Amendment in \textit{Fitzpatrick v. Bitzer}, 427 U.S. 445 (1976), and the Interstate Commerce Clause in \textit{Pennsylvania v. Union Gas Co.}, 491 U.S. 1 (1989). \textit{Fitzpatrick} rested on the decisive chronological and structural fact that the Fourteenth Amendment was adopted \textit{after} the Eleventh Amendment and expressly restricted state sovereignty while empowering Congress to enforce those limits under Section 5 (\href{https://www.law.cornell.edu/supct/html/94-12.ZO.html}{Cornell LII}). By contrast, \textit{Union Gas} was a fractured plurality decision (with Justice White concurring in the judgment without joining Justice Brennan's reasoning) that allowed no principled distinction between the Interstate Commerce Clause and the even more plenary Indian Commerce Clause. Invoking \textit{stare decisis} flexibility in constitutional cases, Rehnquist overruled \textit{Union Gas}: Article I powers antedate the Eleventh Amendment and cannot be used to expand federal judicial power beyond Article III's constitutional limits (\href{https://www.law.cornell.edu/supct/html/94-12.ZO.html}{Cornell LII}). Rehnquist dismissed Justice Souter's textualist reading of the Eleventh Amendment as a ``straw man'' exhibiting ``blind reliance upon the text,'' insisting that \textit{Hans} and \textit{Principality of Monaco v. Mississippi}, 292 U.S. 313 (1934), established state sovereign immunity as an enduring constitutional postulate (\href{https://www.law.cornell.edu/supct/html/94-12.ZO.html}{Cornell LII}).

Finally, on \textit{Ex parte Young}, Rehnquist analogized to \textit{Schweiker v. Chilicky}, 487 U.S. 412 (1988), reasoning that where Congress creates a carefully crafted, limited statutory remedial scheme (\S 2710(d)(7)'s modest mediation and Secretarial-procedures ladder), courts should not permit an \textit{Ex parte Young} suit against a state officer that would expose the officer to full judicial contempt powers and render the statutory scheme superfluous (\href{https://www.law.cornell.edu/supct/html/94-12.ZO.html}{Cornell LII}).

\subsection{Dissenting Opinion --- Justice Stevens}

Justice Stevens argued that the majority committed a ``sharp break with the past'' and an ``affront to a coequal branch'' by stripping Congress of power to provide a federal forum for private enforcement of federal statutes---warning that the ruling would leave victims of state violations of federal copyright, patent, bankruptcy, antitrust, and environmental laws (areas of exclusive federal jurisdiction) without a direct damages remedy against States (\href{https://www.law.cornell.edu/supct/html/94-12.ZD.html}{Cornell LII}).

Stevens contended that neither Justice Iredell's dissent in \textit{Chisholm v. Georgia}, 2 Dall. 419 (1793), nor the Court's decision in \textit{Hans v. Louisiana} constitutionalized state sovereign immunity against congressional abrogation. Rather, both rested on statutory interpretation of the Judiciary Acts of 1789 and 1875, treating sovereign immunity as a rebuttable common-law presumption that Congress could displace by clear legislation (\href{https://www.law.cornell.edu/supct/html/94-12.ZD.html}{Cornell LII}). Furthermore, Stevens argued that because the Eleventh Amendment's plain text applies only to suits by out-of-state or foreign citizens, it has ``absolutely nothing to do'' with federal-question suits brought by a State's own citizens or Indian tribes; treating \textit{Fitzpatrick}'s Section 5 rationale as an exclusive exception was illogical because Congress needs no special post-1795 amendment to displace a nonconstitutional common-law defense (\href{https://www.law.cornell.edu/supct/html/94-12.ZD.html}{Cornell LII}). Stevens also criticized the monarchic, judge-made roots of sovereign immunity (``the King can do no wrong'') as antithetical to American popular sovereignty.

\subsection{Dissenting Opinion --- Justice Souter (Joined by Justices Ginsburg and Breyer)}

In a comprehensive 85-page historical and doctrinal dissent, Justice Souter argued that the Court had, ``for the first time since the founding of the Republic,'' denied Congress authority to subject a State to federal jurisdiction in a suit asserting a federal right (\href{https://www.law.cornell.edu/supct/html/94-12.ZD1.html}{Cornell LII}). Souter developed three main pillars.

First, relying on text, drafting history (including Congress's rejection of Theodore Sedgwick's broader amendment), and Chief Justice Marshall's opinions in \textit{Cohens v. Virginia} and \textit{Osborn v. Bank of the United States}, Souter argued that the Eleventh Amendment merely repealed Article III's Citizen-State Diversity Clauses when a State is a defendant, leaving federal-question (``arising under'') jurisdiction completely untouched (\href{https://www.law.cornell.edu/supct/html/94-12.ZD1.html}{Cornell LII}).

Second, Souter demonstrated that \textit{Hans v. Louisiana}---born out of post-Reconstruction Southern bond repudiation and enforcement pragmatism---had never held that state sovereign immunity was constitutionally immune from congressional abrogation. Even if \textit{Hans} merited retention under \textit{stare decisis} as a federal common-law default rule, the Framers rejected any wholesale constitutional reception of English common law precisely because common-law doctrines were always subject to legislative alteration. Elevating an unwritten common-law immunity above express Article I powers repeated the error of \textit{Lochner v. New York}, 198 U.S. 45 (1905), and Justice Chase's natural-justice dictum in \textit{Calder v. Bull} (\href{https://www.law.cornell.edu/supct/html/94-12.ZD1.html}{Cornell LII}).

Third, Souter dismantled the majority's displacement of \textit{Ex parte Young}. Calling \textit{Young} ``indispensable to the establishment of constitutional government and the rule of law,'' Souter argued that \textit{Young} is not a supplemental \textit{Bivens}-style damages remedy, but a jurisdictional rule for enforcing paramount federal law prospectively against state officers; Congress's statutory procedures in \S 2710(d)(7) would still govern a \textit{Young} suit (just as statutory exhaustion rules govern habeas corpus suits against state wardens), and constitutional avoidance required reading IGRA to preserve \textit{Young} (\href{https://www.law.cornell.edu/supct/html/94-12.ZD1.html}{Cornell LII}).

\section{Examples: Future Applications}

\begin{enumerate}
\item \textbf{Same-Side (Barred by Sovereign Immunity): Private Patent or Copyright Infringement Damages Suit Against a State.} Congress enacts a statute under the Article I Patent and Copyright Clause (Art. I, \S 8, cl. 8) expressly authorizing patent or copyright owners to sue unconsenting States in federal court for money damages. Under \textit{Seminole Tribe}, Congress cannot abrogate state sovereign immunity using an Article I power, no matter how clear its statutory text; unless the statute can be independently sustained as congruent and proportional enforcement legislation under Section 5 of the Fourteenth Amendment (which the Court later rejected in \textit{Florida Prepaid Postsecondary Education Expense Board v. College Savings Bank}, 527 U.S. 627 (1999), and \textit{Allen v. Cooper}, 589 U.S. 248 (2020)), the suit must be dismissed (\href{http://www.stanfordlawreview.org/wp-content/uploads/sites/3/2010/04/chemerinsky.pdf}{Stanford Law Review}).

\item \textbf{Same-Side (Barred by Sovereign Immunity): Federal Statutory Age- or Disability-Discrimination Damages Suit Exceeding Section 5.} Congress amends a federal employment statute under the Commerce Clause and Section 5 of the Fourteenth Amendment to authorize state employees to sue their State employers for money damages for age discrimination. Under \textit{Seminole Tribe}, the Commerce Clause cannot support abrogation; therefore, the suit survives \textit{only} if the statute is a valid, congruent, and proportional exercise of Section 5 of the Fourteenth Amendment under \textit{City of Boerne v. Flores}. Because age is not a suspect classification and Congress lacked a sufficient record of unconstitutional state age discrimination, the abrogation fails. (The Court applied \textit{Seminole Tribe} in exactly this way in \textit{Kimel v. Florida Board of Regents}, 528 U.S. 62 (2000), and \textit{Board of Trustees of the Univ. of Ala. v. Garrett}, 531 U.S. 356 (2001)) (\href{http://www.stanfordlawreview.org/wp-content/uploads/sites/3/2010/04/chemerinsky.pdf}{Stanford Law Review}).

\item \textbf{Opposite-Side (Permitted): Family-Leave Damages Suit Validly Enacted Under Section 5 of the Fourteenth Amendment.} Congress enacts a federal statute expressly authorizing state employees to sue States for money damages for violating family-care leave protections, backed by an extensive legislative record showing pervasive, unconstitutional gender-based stereotyping by state employers. Because \textit{Seminole Tribe} expressly reaffirmed \textit{Fitzpatrick v. Bitzer}, Congress retains full constitutional power to abrogate state sovereign immunity when enacting congruent and proportional legislation under Section 5 of the Fourteenth Amendment to remedy heightened-scrutiny Equal Protection violations (as the Court held in \textit{Nevada Dept. of Human Resources v. Hibbs}, 538 U.S. 721 (2003)).

\item \textbf{Opposite-Side (Permitted): Traditional \textit{Ex parte Young} Suit Against a State Official Under a General Enforcement Provision.} A telecommunications carrier sues state public utility commissioners in their official capacities seeking prospective injunctive relief to stop them from enforcing a state rate order that violates a federal telecommunications statute containing a general judicial-review provision without a narrowly limited remedial ladder like IGRA's. Under \textit{Seminole Tribe}, \textit{Ex parte Young} remains fully available whenever a plaintiff seeks prospective injunctive relief to end an ongoing violation of federal law and Congress has not enacted a detailed, restrictive remedial scheme signaling intent to displace officer suits (as confirmed in \textit{Verizon Maryland Inc. v. Public Service Comm'n of Maryland}, 535 U.S. 635 (2002)).

\item \textbf{Fence-Sitter: Abrogation Under Non-Commerce Article I Powers Claimed to Involve a ``Plan of the Convention'' Waiver (e.g., Bankruptcy or War Powers).} Suppose Congress authorizes private suits against unconsenting States under the Article I Bankruptcy Clause (Art. I, \S 8, cl. 4) to recover preferential transfers, or under the Article I War Powers to enforce returning servicemembers' reemployment rights. On one hand, \textit{Seminole Tribe}'s sweeping language declared that ``Article I cannot be used to circumvent the constitutional limitations placed upon federal jurisdiction.'' On the other hand, \textit{Seminole Tribe} acknowledged that States lack immunity where there was a ``surrender of this immunity in the plan of the convention.'' This tension makes non-Commerce Article I powers genuine fence-sitters: the Court later held that States surrendered immunity in the plan of the Convention for \textit{in rem} bankruptcy proceedings (\textit{Central Virginia Community College v. Katz}, 546 U.S. 356 (2006)), eminent domain (\textit{PennEast Pipeline Co. v. New Jersey}, 594 U.S. 482 (2021)), and war-powers military reemployment (\textit{Torres v. Texas Dept. of Public Safety}, 597 U.S. 580 (2022)), carving notable exceptions into \textit{Seminole Tribe}'s seemingly categorical Article I bar (\href{https://texaslawreview.org/state-sovereign-immunity-after-the-revolution/}{Texas Law Review}).
\end{enumerate}

\section{Critique}

\textbf{Originalist and Textualist Perspectives:} \textit{Seminole Tribe} has long been a flashpoint among originalists and textualists. Critics note an acute methodological irony: Justices Scalia and Thomas---champions of strict textualism---joined an opinion that expressly disparaged ``blind reliance upon the text of the Eleventh Amendment'' in favor of an unwritten, atextual ``presupposition'' of state sovereign immunity (\href{https://www.law.cornell.edu/supct/html/94-12.ZO.html}{Cornell LII}). Strict textualists and original-public-meaning scholars (echoing Justice Souter's dissent and earlier work by Akhil Amar and William Fletcher) argue that the Eleventh Amendment's text bars only diversity suits by out-of-state or foreign citizens, and that nothing in Article III or the Eleventh Amendment strips Congress of power to abrogate common-law immunity in federal-question cases. More recently, however, a prominent school of ``original-law originalists''---notably William Baude and Stephen E. Sachs---has offered a sophisticated defense of \textit{Seminole Tribe}'s bottom line (\href{https://pennlawreview.com/2021/05/11/the-misunderstood-eleventh-amendment/}{University of Pennsylvania Law Review}; \href{https://www.scotusblog.com/2016/01/academic-highlight-baude-and-sachs-on-state-sovereign-immunity/}{SCOTUSblog}). Under their ``constitutional backdrop'' theory, state sovereign immunity was a preexisting general-law rule preserved at the Founding that neither Article III nor Article I authorized Congress to abrogate, while the Eleventh Amendment simply prevented federal courts from misconstruing Article III to override that backdrop.

\textbf{Progressive and Doctrinal Perspectives:} Progressive and federal-courts scholars---including Vicki C. Jackson (\href{https://www.nyulawreview.org/wp-content/uploads/2018/08/NYULawReview-72-3-Jackson.pdf}{NYU Law Review}) and Erwin Chemerinsky (\href{http://www.stanfordlawreview.org/wp-content/uploads/sites/3/2010/04/chemerinsky.pdf}{Stanford Law Review})---heavily criticize \textit{Seminole Tribe} on both doctrinal and practical grounds. First, by foreclosing Article I abrogation, \textit{Seminole Tribe} forced Congress to funnel nearly all private enforcement remedies against States through Section 5 of the Fourteenth Amendment---precisely as the Rehnquist Court simultaneously constricted Section 5 through the ``congruence and proportionality'' test of \textit{City of Boerne v. Flores}, creating a federalism pincer that immunized States from damages under federal civil-rights, labor, and intellectual-property laws (\href{http://www.stanfordlawreview.org/wp-content/uploads/sites/3/2010/04/chemerinsky.pdf}{Stanford Law Review}). Second, Vicki Jackson warned that \textit{Seminole Tribe}'s curtailment of \textit{Ex parte Young} threatened to eviscerate the primary mechanism for ensuring state compliance with supreme federal law whenever Congress includes structured enforcement provisions (\href{https://www.nyulawreview.org/wp-content/uploads/2018/08/NYULawReview-72-3-Jackson.pdf}{NYU Law Review}). Third, critics point out that relying on federal-agency enforcement suits (rather than private citizen suits) is unrealistic given limited federal enforcement budgets, leaving federal rights without effective remedies.

\textbf{Underlying Competing Values:} \textit{Seminole Tribe} crystallizes the enduring clash between \textit{state sovereignty (dignity and fiscal autonomy in a dual-sovereign federal system)} and \textit{federal supremacy and the rule-of-law principle that every federal right deserves an effective judicial remedy}. The majority prioritizes protecting States from the coercive indignity and fiscal burdens of private federal litigation unless they consent or Congress enforces the post-Civil War Fourteenth Amendment. The dissents prioritize democratic accountability through Congress's plenary Article I powers and the Supremacy Clause imperative that States be judicially answerable when they violate valid federal law.

\section{Key Quotations}

\begin{quote}
``Even when the Constitution vests in Congress complete lawmaking authority over a particular area, the Eleventh Amendment prevents congressional authorization of suits by private parties against unconsenting States. The Eleventh Amendment restricts the judicial power under Article III, and Article I cannot be used to circumvent the constitutional limitations placed upon federal jurisdiction.''
\end{quote}
\noindent --- Chief Justice Rehnquist, opinion of the Court (\href{https://www.law.cornell.edu/supct/html/94-12.ZO.html}{Cornell LII}; \href{https://www.courtlistener.com/opinion/118011/seminole-tribe-of-florida-v-florida/}{CourtListener}).

\begin{quote}
``The dissent's lengthy analysis of the text of the Eleventh Amendment is directed at a straw man---we long have recognized that blind reliance upon the text of the Eleventh Amendment is `` `to strain the Constitution and the law to a construction never imagined or dreamed of.' '' \textit{Monaco}, 292 U.S. at 326, quoting \textit{Hans}, 134 U.S. at 15.''
\end{quote}
\noindent --- Chief Justice Rehnquist, opinion of the Court (\href{https://www.law.cornell.edu/supct/html/94-12.ZO.html}{Cornell LII}; \href{https://www.courtlistener.com/opinion/118011/seminole-tribe-of-florida-v-florida/}{CourtListener}).

\begin{quote}
``[T]herefore, where Congress has prescribed a detailed remedial scheme for the enforcement against a State of a statutorily created right, a court should hesitate before casting aside those limitations and permitting an action against a state officer based upon \textit{Ex parte Young}.''
\end{quote}
\noindent --- Chief Justice Rehnquist, opinion of the Court (\href{https://www.law.cornell.edu/supct/html/94-12.ZO.html}{Cornell LII}; \href{https://www.courtlistener.com/opinion/118011/seminole-tribe-of-florida-v-florida/}{CourtListener}).

\begin{quote}
``The majority's opinion does not simply preclude Congress from establishing the rather curious statutory scheme under which Indian tribes may seek the aid of a federal court to secure a State's good faith negotiations over gaming regulations. Rather, it prevents Congress from providing a federal forum for a broad range of actions against States, from those sounding in copyright and patent law, to those concerning bankruptcy, environmental law, and the regulation of our vast national economy.''
\end{quote}
\noindent --- Justice Stevens, dissenting (\href{https://www.law.cornell.edu/supct/html/94-12.ZD.html}{Cornell LII}; \href{https://www.courtlistener.com/opinion/118011/seminole-tribe-of-florida-v-florida/}{CourtListener}).

\begin{quote}
``In holding the State of Florida immune to suit under the Indian Gaming Regulatory Act, the Court today holds for the first time since the founding of the Republic that Congress has no authority to subject a State to the jurisdiction of a federal court at the behest of an individual asserting a federal right.''
\end{quote}
\noindent --- Justice Souter, dissenting (\href{https://www.law.cornell.edu/supct/html/94-12.ZD1.html}{Cornell LII}; \href{https://www.courtlistener.com/opinion/118011/seminole-tribe-of-florida-v-florida/}{CourtListener}).


\end{document}
